% --- the Chinese edition --------------------------------------------------
% Loaded after preamble.tex and after report-zh-strings.tex, which is
% generated by scripts/translate.mjs and renews the document's own words.
%
% This file is the whole difference between the two editions apart from those
% words. It is deliberately short: the layout, the palette, the page, the
% running heads, the callouts and the title page are written once, in
% preamble.tex, and are not repeated here.
%
% Four things genuinely differ.

% --- 1. the faces --------------------------------------------------------
% preamble.tex sets ONE Chinese cut, at Regular, because the English edition
% only ever needed han glyphs for a statutory quotation. A Chinese document
% sets its whole body in that face, and the moment it does, the missing bold
% matters: fontspec with no BoldFont declared silently leaves \bfseries at the
% regular weight, so every **emphasis** in the report — and every callout
% title, and every table header — would come out looking like body text, with
% no error anywhere. That is the same class of failure as a dropped glyph, and
% it is fixed the same way: fonts.py instances a second cut and it is declared
% here.
\setCJKmainfont{NotoSerifSC-report.otf}[Path=fonts/,
  BoldFont = NotoSerifSC-report-Bold.otf]
\setCJKmonofont{NotoSerifSC-report.otf}[Path=fonts/,
  BoldFont = NotoSerifSC-report-Bold.otf]
% Headings are NOT given a heavier Chinese cut, on purpose. The Latin display
% face in this document is Fraunces at wght 400 — \titleformat asks for
% \displayfont\LARGE, not for bold — so the Chinese main face at Regular is the
% weight that matches it. A bold Chinese heading beside a regular Latin one is
% the mismatch, not the fix.

% --- 2. the leading ------------------------------------------------------
% Set in the front matter (linestretch 1.45, against the English edition's
% 1.15) because that is where Quarto will accept it. A page of han characters
% carries far more ink per line than a page of Latin at the same size: at 1.15
% the lines lock together and the eye loses its place returning to the left
% margin.
%
% Paragraph indentation is NOT changed. Chinese convention indents the first
% line by two characters and sets no space between paragraphs; this document
% separates paragraphs by space, in both editions, because the two are meant to
% be readable side by side, page against page. A reader comparing them should
% find the same document, not two different ones.
%
% What IS needed is xeCJK's widow check: without it a single han character can
% be left alone on the last line of a paragraph, which in Chinese reads as
% damage rather than as a short line.
\xeCJKsetup{CheckSingle=true}
% preamble.tex hands the horizontal ellipsis to the Latin face, because in an
% English sentence that is where it belongs. In Chinese it is Chinese
% punctuation — full-width, and set as a pair — so it comes back to the CJK
% face here. fonts.py keeps it in the subset for both editions.
\xeCJKDeclareCharClass{CJK}{"2026}

% --- 3. the names LaTeX supplies ----------------------------------------
% Quarto's own answer to a Chinese document is to load babel or polyglossia for
% it. Neither sets Chinese without a locale package, and both then argue with
% xeCJK under tectonic — so the .qmd keeps `lang: en` and the names are set
% here, explicitly, from the project's lexicon.
%
% \AfterEndPreamble, not \renewcommand and not \AtBeginDocument. Pandoc's
% template writes these six names itself, in English, inside an
% \AtBeginDocument block of its own that sits AFTER this file — so a plain
% \renewcommand here is overwritten before the first page, and an
% \AtBeginDocument here runs before the template's because hooks fire in the
% order they were declared. The contents page read "Table of contents" in the
% Chinese edition until this was found. etoolbox's \AfterEndPreamble runs after
% every \AtBeginDocument hook, which is the only place that wins.
\AfterEndPreamble{%
  \renewcommand{\contentsname}{\sccontents}%
  \renewcommand{\listfigurename}{\sclistOfFigures}%
  \renewcommand{\listtablename}{\sclistOfTables}%
  \renewcommand{\figurename}{\scfigure}%
  \renewcommand{\tablename}{\sctable}%
  \renewcommand{\indexname}{\scindex}%
  % The document's real language, for anything reading the file rather than the
  % page: a screen reader, a search index, a print shop's preflight.
  \hypersetup{pdflang=zh-CN}}

% --- 4. the untranslated-passage mark -----------------------------------
% A small superscript EN in the warning colour, at the head of any passage
% still in English. Inline rather than an environment, because the same mark
% has to work at the head of a paragraph, inside a table cell, in a figure
% caption and in a heading — an environment works in only the first of those.
%
% \texorpdfstring keeps it out of the PDF bookmarks, where hyperref would
% otherwise write the macro's expansion into a plain-text string. The count is
% on the title page, so the reader is told how much of the document this
% applies to before meeting the first instance of it.
%
% It is also suppressed in the three lists at the front. A heading's mark is
% written into the .toc with the heading, and a contents page with an EN beside
% ninety of its lines stops being a map of the document — which is the one job
% a contents page has. The flag is read where the mark is TYPESET, so switching
% it off around each list is enough; the heading itself, and the running head
% above it, still carry the mark.
\newif\ifzhmarkon \zhmarkontrue
\DeclareRobustCommand{\zhmark}{%
  \texorpdfstring{\ifzhmarkon\textsuperscript{\uifont\tiny\color{warn}EN}\fi}{}}
\let\sczhtoc\tableofcontents \def\tableofcontents{\zhmarkonfalse\sczhtoc\zhmarkontrue}
\let\sczhlof\listoffigures   \def\listoffigures{\zhmarkonfalse\sczhlof\zhmarkontrue}
\let\sczhlot\listoftables    \def\listoftables{\zhmarkonfalse\sczhlot\zhmarkontrue}
